\section{Application: an EVM rollup calldata commitment}
\label{sec:rollup-eth}

FlatMerkle applies equally to an EVM-compatible rollup's data-availability
problem, though for a narrower reason than on Bitcoin: on an EVM chain,
FlatMerkle is not the cheapest way to publish a byte, only the cheapest way
we find to make an already-published, already-expensive byte cheaply
disputable while keeping it \emph{persistently} available. We work this out
in full for Ethereum first, then redo the same four options against
Rootstock \citep{rsk2022}, a Bitcoin-merge-mined, EVM-compatible sidechain,
and against Arbitrum \citep{arbos40}, an Ethereum rollup in its own right --
not as further illustrations of the same result, but because the three
chains turn out to disagree with each other on which option wins, for a
reason worth understanding: the ranking depends on specific,
independently adoptable protocol rules, not on anything intrinsic to
FlatMerkle or to the EVM itself. The full, from-scratch derivation --
every gas constant cited to its defining EIP, RSKIP, or ArbOS release
note, every non-protocol approximation flagged as such -- is given in
Appendix~\ref{app:eth-da}; this section states the setup, the four
options being compared, and (\S\ref{sec:rollup-eth-results}) the result.

\subsection{Why calldata, not blobs}
\label{sec:rollup-eth-setup}

A rollup sequencer can publish its batch data to Ethereum in ordinary
transaction \emph{calldata}, costed at 16 gas per non-zero byte (4 gas per
zero byte), a protocol constant unchanged since 2019
\citep{eip2028}. Since Dencun, it can instead publish the same bytes in a
\emph{blob} \citep{eip4844}: a separate, independently fee-market-priced
resource (blob gas is charged 1:1 against a blob's $131{,}072$ bytes,
BLOB\_GAS\_PER\_BLOB $=2^{17}$), historically far cheaper per byte than
calldata when blobspace demand is low -- though this ratio floats with two
independent, unrelated fee markets and is not a fixed protocol constant, so
we do not rely on any particular multiple here. The catch, stated directly
in EIP-4844's own text, is that blobs are not meant to be kept: ``[the
consensus layer] blobs must be kept for at least a certain period[,] around
18 days'' (\texttt{MIN\_EPOCHS\_FOR\_BLOB\_SIDECARS\_REQUESTS}
$=4096$ epochs) \citep{eip4844}, after which nodes are free to discard
them. A rollup that wants Ethereum itself to remain a permanent record --
so that a new L2 node can reconstruct the entire rollup state by replaying
L1 history alone, indefinitely, the way this paper's other applications
assume the underlying chain's data stays available -- cannot rely on
blobs for that role; ordinary calldata, once included in a block, has no
comparable expiry. This section is about that persistent, expensive-by-design
case: given a sequencer already pays calldata's price to publish $P$ bytes
permanently, what is the cheapest way to also get a compact, on-chain,
disputable commitment to those bytes? On Rootstock the question is not even
a choice: there is no blob-carrying transaction type there at all, so
calldata (or contract storage) is the only way to publish data on-chain in
the first place, making a cheap on-chain commitment mechanism, if anything,
more directly load-bearing there than on Ethereum, not less.

\subsection{Four ways to commit to published calldata}
\label{sec:rollup-eth-options}

\begin{enumerate}
\item[A.] Copy the calldata into contract storage (\texttt{SSTORE}) and
  keep it there.
\item[B.] Build an ordinary Keccak256 Merkle tree over the calldata inside
  the contract, store only the 32-byte root, and discard everything else.
\item[C.] Build a FlatMerkle structured file over the data \emph{off-chain},
  submit the whole structured file as calldata, and hash it on-chain in one
  call to the SHA-256 precompile (address \texttt{0x02}), storing only the
  digest.
\item[D.] Build the FlatMerkle structured file \emph{inside} the contract,
  from ordinary (unstructured) calldata, hashing it incrementally via the
  BLAKE2 compression-function precompile (address \texttt{0x09},
  \citep{eip152}), storing only the final digest.
\end{enumerate}

Option D uses BLAKE2b's compression function $F(h, m, t, f)$ as $\C$. Besides
the chaining value $h$ and the message block $m$, $F$ takes a byte counter
$t$ and a final-block flag $f$, and the precompile accepts any values for
both. A FlatMerkle instantiation must fix them from the item's position
alone -- for example, $t$ the running byte count $128 \cdot (\text{item
index} + 1)$ and $f = 0$ for every item, with any finalization applied as a
separate fixed step -- in the contract that builds the file and in any verifier that
later checks a proof against its digest, and must never take them from a
proof. Otherwise a proof could vary them, and
Assumption~\ref{ass:freestart}, which is stated for a fixed function of
$(h,m)$, would no longer cover the verifier's computation
(Remark~\ref{rem:freestart-history} discusses BLAKE2b's own security
margin).

Options B and D need no fraud-proof machinery of the kind
Section~\ref{sec:fraud-proofs} gives for the Bitcoin annex application: the
commitment is built by trusted contract code from data the contract itself
received, so it cannot be non-canonical the way an externally supplied
FlatMerkle file (option C, or the Bitcoin annex of
Section~\ref{sec:bitcoin-annex}) can be. A later dispute over the
\emph{data itself} (was this the right previous state? did this transaction
really produce this result?) is a separate, application-level question this
section does not address, matching Section~\ref{sec:rollup-da}'s treatment
of the analogous Bitcoin case.

\subsection{Result}
\label{sec:rollup-eth-results}

\begin{table}[htbp]
\centering
\begin{tabular}{@{}lrrrr@{}}
\toprule
Option & Calldata & Ethereum & Rootstock & Arbitrum \\
\midrule
A. SSTORE, persistent      & $P$  & $706.625$ & $641$            & $706.625$ \\
B. Keccak tree, discarded  & $P$  & $40$      & $17.3$--$18.4$   & $17.3$--$18.4$ \\
C. FlatMerkle, off-chain   & $2P$ & $80$      & $32.75$--$43.5$  & $32.75$--$34.1$ \\
D. FlatMerkle, on-chain    & $P$  & $40$      & $\approx\!27.5$  & $16.2$--$18.1$ \\
\bottomrule
\end{tabular}
\caption{Total transaction gas per original byte $P$ (asymptotic),
current protocol constants on three chains, derived in
Appendix~\ref{app:eth-da}. Ethereum: EIP-7623 calldata floor and EIP-2929
warm-call discount both present (floor dominates). Rootstock: neither.
Arbitrum: EIP-2929 discount, no EIP-7623 floor -- the one combination
where D is not a clear loser against B.}
\label{tab:eth-da-costs}
\end{table}

The three chains disagree with each other, and the disagreement is the
result. On Ethereum, option D does not beat option B: they tie, both at
the EIP-7623 floor, which is indifferent to a more than tenfold spread in
the two options' real hashing cost. On Rootstock, option D loses to
option B by a real margin ($\approx\!27.5$ against
$\approx\!17.3$--$18.4$ gas/byte): with no floor to hide behind, raw
efficiency matters again, and Rootstock's missing warm-call discount
taxes option D's many separate BLAKE2F calls more than it costs option
B's zero calls. On Arbitrum -- no floor, but \emph{with} the warm-call
discount -- option D is, for the first time, at least competitive with
option B and plausibly a little cheaper, though not decisively so given
the uncertainty in our own ``glue code'' estimate (Appendix~\ref{app:eth-da}).
So the premise this section started from -- that FlatMerkle built inside
the contract should be strictly cheaper than a plain Merkle tree -- is
only borne out, and only weakly, on the one chain of the three that
combines both relevant protocol properties (a missing calldata floor and
a present warm-call discount); it does not hold in general, and which of
B or D wins is a fact about a given chain's own protocol rules, not about
FlatMerkle or about ``the EVM'' as a single target. FlatMerkle-in-contract
is, however, a large and consistent improvement over building the
structured file off-chain (option C) and over persisting the raw data in
storage (option A) on all three chains -- roughly $1.5$--$2\times$
cheaper than C and well over an order of magnitude cheaper than A
throughout, for reasons and with a fuller discussion of what this
does and does not establish given in Appendix~\ref{app:eth-da}.
Ethereum's pending Glamsterdam fork, as currently drafted, would raise the
calldata floor to $64$ gas per byte; that raises the Ethereum column for
B, C and D to $64$, $128$ and $64$ but leaves the ranking unchanged
(Appendix~\ref{app:eth-da}).
